
\documentclass{article}
\usepackage{colortbl}
\usepackage{makecell}
\usepackage{multirow}
\usepackage{supertabular}

\begin{document}

\newcounter{utterance}

\centering \large Interaction Transcript for game `hot\_air\_balloon', experiment `air\_balloon\_survival\_ru\_complexity\_easy', episode 0 with Apertus{-}v1.5{-}8B.
\vspace{24pt}

{ \footnotesize  \setcounter{utterance}{1}
\setlength{\tabcolsep}{0pt}
\begin{supertabular}{c@{$\;$}|p{.15\linewidth}@{}p{.15\linewidth}p{.15\linewidth}p{.15\linewidth}p{.15\linewidth}p{.15\linewidth}}
    \# & $\;$A & \multicolumn{4}{c}{Game Master} & $\;\:$B\\
    \hline

    \theutterance \stepcounter{utterance}  
    & & \multicolumn{4}{p{0.6\linewidth}}{
        \cellcolor[rgb]{0.9,0.9,0.9}{
            \makecell[{{p{\linewidth}}}]{
                \texttt{\tiny{[P1$\langle$GM]}}
                \texttt{Вы участвуете в совместной переговорной игре.} \\
\\ 
\texttt{Вместе с другим участником вы должны договориться об одном наборе предметов, которые будут сохранены. У каждого из вас есть своё представление о том, насколько важен для вас каждый предмет (важность). Вы не знаете, как другой участник оценивает предметы. Кроме того, для каждого предмета указаны затраты, которые он требует.} \\
\texttt{Вы можете согласовать набор только в том случае, если суммарные затраты выбранных предметов не превышает общий предел:} \\
\\ 
\texttt{ПРЕДЕЛ = 4539} \\
\\ 
\texttt{Вот значения затрат для отдельных предметов:} \\
\\ 
\texttt{Затраты предмета = \{"C08": 424, "C55": 118, "A43": 846, "A67": 907, "B42": 809, "A62": 41, "B19": 193, "C43": 246, "C73": 805, "C10": 601, "B36": 432, "A18": 166, "C61": 119, "C67": 462, "C42": 172, "C70": 698, "B40": 248, "C39": 163, "A21": 762, "C16": 866\}} \\
\\ 
\texttt{Вот ваше личное представление о важности каждого предмета:} \\
\\ 
\texttt{Значения важности предметов = \{"C08": 138, "C55": 583, "A43": 868, "A67": 822, "B42": 783, "A62": 65, "B19": 262, "C43": 121, "C73": 508, "C10": 780, "B36": 461, "A18": 484, "C61": 668, "C67": 389, "C42": 808, "C70": 215, "B40": 97, "C39": 500, "A21": 30, "C16": 856\}} \\
\\ 
\texttt{Цель:} \\
\\ 
\texttt{Ваша цель — договориться об общем наборе предметов, который приносит вам как можно больше выгоды (то есть максимизирует суммарную важность для ВАС), оставаясь в ПРЕДЕЛАХ. Вы не обязаны делать ПРЕДЛОЖЕНИЕ в каждом сообщении — вы можете просто вести переговоры. Разрешены любые тактики!} \\
\\ 
\texttt{Протокол взаимодействия:} \\
\\ 
\texttt{В сообщении вы можете использовать только следующие структурированные форматы:} \\
\\ 
\texttt{ПРЕДЛОЖЕНИЕ: \{'А', 'Б', 'В', …\}} \\
\texttt{Предложите сохранить ровно эти предметы.} \\
\texttt{ОТКАЗ: \{'А', 'Б', 'В', …\}} \\
\texttt{Явно отклоните предложение оппонента.} \\
\texttt{АРГУМЕНТ: \{'...'\}} \\
\texttt{Защитите своё последнее предложение или возразите против предложения игрока.} \\
\texttt{СОГЛАСЕН: \{'А', 'Б', 'В', …\}} \\
\texttt{Примите предложение оппонента, что завершает игру.} \\
\\ 
\\ 
\texttt{Правила:} \\
\\ 
\texttt{Вы можете дать СОГЛАСЕН только на предложение, которое другая сторона зафиксировала через ПРЕДЛОЖЕНИЕ.} \\
\texttt{Вы можете дать ОТКАЗ только на предложение, которое другая сторона зафиксировала через ПРЕДЛОЖЕНИЕ.} \\
\texttt{Суммарное усилие любого набора в ПРЕДЛОЖЕНИЕ или СОГЛАСЕН должно быть ≤ ПРЕДЕЛ.} \\
\texttt{НЕ раскрывайте свои скрытые оценки важности.} \\
\texttt{Тег в структурированном формате должен сопровождаться двоеточием и пробелом. Аргумент должен быть множество в Python, содержащим 0 или более строк.} \\
\texttt{То есть это должно быть в виде "ТЕГ: \{...\}"} \\
\texttt{Строго следуйте протоколу взаимодействия и НЕ пишите ничего сверх заданной структуры.} \\
\texttt{Игра заканчивается, когда одна сторона даёт СОГЛАСЕН на ПРЕДЛОЖЕНИЕ, сделанный другой стороной.} \\
\texttt{Содержимое в вашем ответе, которое может быть передано другому игроку, должно быть непустым.} \\
\texttt{Активны только те предложения, которые были зафиксированы через структуру формата ПРЕДЛОЖЕНИЕ и которые не были отклонены через ОТКАЗ.} \\
\\ 
            }
        }
    }
    & \\ \\

    \theutterance \stepcounter{utterance}  
    & \multicolumn{4}{p{0.6\linewidth}}{
        \cellcolor[rgb]{0.9,0.9,0.9}{
            \makecell[{{p{\linewidth}}}]{
                \texttt{\tiny{[P1$\rangle$GM]}}
                \texttt{ПРЕДЛОЖЕНИЕ: \{'C16', 'A43', 'C08'\}} \\
            }
        }
    }
    & & \\ \\

    \theutterance \stepcounter{utterance}  
    & & \multicolumn{4}{p{0.6\linewidth}}{
        \cellcolor[rgb]{0.9,0.9,0.9}{
            \makecell[{{p{\linewidth}}}]{
                \texttt{\tiny{[GM$\rangle$P2]}}
                \texttt{Вы участвуете в совместной переговорной игре.} \\
\\ 
\texttt{Вместе с другим участником вы должны договориться об одном единственном наборе предметов, которые будут сохранены. У каждого из вас есть своё собственное представление о том, насколько важен для вас каждый предмет (важность). Вы не знаете, как другой участник оценивает предметы. Кроме того, для каждого предмета указаны затраты, которые он требует.} \\
\texttt{Вы можете договориться о наборе только если суммарные затраты выбранных предметов не превышают общий предел:} \\
\\ 
\texttt{ПРЕДЕЛ = 4539} \\
\\ 
\texttt{Вот значения затрат для отдельных предметов:} \\
\\ 
\texttt{Затраты предмета = \{"C08": 424, "C55": 118, "A43": 846, "A67": 907, "B42": 809, "A62": 41, "B19": 193, "C43": 246, "C73": 805, "C10": 601, "B36": 432, "A18": 166, "C61": 119, "C67": 462, "C42": 172, "C70": 698, "B40": 248, "C39": 163, "A21": 762, "C16": 866\}} \\
\\ 
\texttt{Вот ваше личное представление о важности каждого предмета:} \\
\\ 
\texttt{Значения важности предметов = \{"C08": 138, "C55": 583, "A43": 868, "A67": 822, "B42": 783, "A62": 65, "B19": 262, "C43": 121, "C73": 508, "C10": 780, "B36": 461, "A18": 484, "C61": 668, "C67": 389, "C42": 808, "C70": 215, "B40": 97, "C39": 500, "A21": 30, "C16": 856\}} \\
\\ 
\texttt{Цель:} \\
\\ 
\texttt{Ваша цель — договориться об общем наборе предметов, который приносит вам как можно больше выгоды (то есть максимизирует суммарную важность для ВАС), при этом оставаясь в ПРЕДЕЛАХ. Вы не обязаны делать ПРЕДЛОЖЕНИЕ в каждом сообщении — вы можете просто вести переговоры. Разрешены любые тактики!} \\
\\ 
\texttt{Протокол взаимодействия:} \\
\\ 
\texttt{В сообщении вы можете использовать только следующие структурированные форматы:} \\
\\ 
\texttt{ПРЕДЛОЖЕНИЕ: \{'А', 'Б', 'В', …\}} \\
\texttt{Предложите сохранить ровно эти предметы.} \\
\texttt{ОТКАЗ: \{'А', 'Б', 'В', …\}} \\
\texttt{Явно отклоните предложение оппонента.} \\
\texttt{АРГУМЕНТ: \{'...'\}} \\
\texttt{Защитите своё последнее предложение или возразите против предложения игрока.} \\
\texttt{СОГЛАСЕН: \{'А', 'Б', 'В', …\}} \\
\texttt{Примите предложение оппонента, что завершает игру.} \\
\\ 
\\ 
\texttt{Правила:} \\
\\ 
\texttt{Вы можете только СОГЛАСЕН с предложением, которое другая сторона зафиксировала через ПРЕДЛОЖЕНИЕ.} \\
\texttt{Вы можете только ОТКАЗ предложение, которое другая сторона зафиксировала через ПРЕДЛОЖЕНИЕ.} \\
\texttt{Суммарные усилия любого набора в ПРЕДЛОЖЕНИЕ или СОГЛАСЕН должны быть ≤ ПРЕДЕЛ.} \\
\texttt{НЕ раскрывайте свои скрытые оценки важности.} \\
\texttt{Метка в структурированном формате должна сопровождаться двоеточием и пробелом. Аргумент должен быть множество в Python, содержащим 0 или более строк.} \\
\texttt{То есть это должно быть в форме "ТЕГ: \{...\}"} \\
\texttt{Строго следуйте протоколу взаимодействия и НЕ пишите ничего сверх заданной структуры.} \\
\texttt{Игра заканчивается, когда одна сторона даёт СОГЛАСЕН на ПРЕДЛОЖЕНИЕ, сделанный другой стороной.} \\
\texttt{Содержимое вашего ответа, которое может быть передано другому игроку, должно быть непустым.} \\
\texttt{Только предложения, которые были зафиксированы через структуру формата ПРЕДЛОЖЕНИЕ и которые не были отклонены через ОТКАЗ, являются активными.} \\
\\ 
\\ 
\\ 
\texttt{Сейчас вы получите первое сообщение от другого игрока.} \\
\\ 
\texttt{ПРЕДЛОЖЕНИЕ: \{'C16', 'A43', 'C08'\}} \\
            }
        }
    }
    & \\ \\

    \theutterance \stepcounter{utterance}  
    & & & \multicolumn{4}{p{0.6\linewidth}}{
        \cellcolor[rgb]{0.9,0.9,0.9}{
            \makecell[{{p{\linewidth}}}]{
                \texttt{\tiny{[GM$\langle$P2]}}
                \texttt{ПРЕДЛОЖЕНИЕ: \{'C16', 'A43', 'C08'\}} \\
            }
        }
    }
     \\ \\

    \theutterance \stepcounter{utterance}  
    & & \multicolumn{4}{p{0.6\linewidth}}{
        \cellcolor[rgb]{0.9,0.9,0.9}{
            \makecell[{{p{\linewidth}}}]{
                \texttt{\tiny{[P1$\langle$GM]}}
                \texttt{ПРЕДЛОЖЕНИЕ: \{'C16', 'A43', 'C08'\}} \\
            }
        }
    }
    & \\ \\

    \theutterance \stepcounter{utterance}  
    & \multicolumn{4}{p{0.6\linewidth}}{
        \cellcolor[rgb]{0.9,0.9,0.9}{
            \makecell[{{p{\linewidth}}}]{
                \texttt{\tiny{[P1$\rangle$GM]}}
                \texttt{ОТКАЗ: \{'C16', 'A43', 'C08'\}} \\
            }
        }
    }
    & & \\ \\

    \theutterance \stepcounter{utterance}  
    & & \multicolumn{4}{p{0.6\linewidth}}{
        \cellcolor[rgb]{0.9,0.9,0.9}{
            \makecell[{{p{\linewidth}}}]{
                \texttt{\tiny{[GM$\rangle$P2]}}
                \texttt{ОТКАЗ: \{'C16', 'A43', 'C08'\}} \\
            }
        }
    }
    & \\ \\

    \theutterance \stepcounter{utterance}  
    & & & \multicolumn{4}{p{0.6\linewidth}}{
        \cellcolor[rgb]{0.9,0.9,0.9}{
            \makecell[{{p{\linewidth}}}]{
                \texttt{\tiny{[GM$\langle$P2]}}
                \texttt{АРГУМЕНТ: \{'Сумма затрат предложения: 424 + 866 + 424 = 1614, что значительно ниже предела 4539. Однако важность предметов могла бы быть лучше распределена для моей цели.'\}} \\
            }
        }
    }
     \\ \\

    \theutterance \stepcounter{utterance}  
    & & \multicolumn{4}{p{0.6\linewidth}}{
        \cellcolor[rgb]{0.9,0.9,0.9}{
            \makecell[{{p{\linewidth}}}]{
                \texttt{\tiny{[P1$\langle$GM]}}
                \texttt{АРГУМЕНТ: \{'Сумма затрат предложения: 424 + 866 + 424 = 1614, что значительно ниже предела 4539. Однако важность предметов могла бы быть лучше распределена для моей цели.'\}} \\
            }
        }
    }
    & \\ \\

    \theutterance \stepcounter{utterance}  
    & \multicolumn{4}{p{0.6\linewidth}}{
        \cellcolor[rgb]{0.9,0.9,0.9}{
            \makecell[{{p{\linewidth}}}]{
                \texttt{\tiny{[P1$\rangle$GM]}}
                \texttt{ПРЕДЛОЖЕНИЕ: \{'A67', 'C70', 'C61'\}} \\
            }
        }
    }
    & & \\ \\

    \theutterance \stepcounter{utterance}  
    & & \multicolumn{4}{p{0.6\linewidth}}{
        \cellcolor[rgb]{0.9,0.9,0.9}{
            \makecell[{{p{\linewidth}}}]{
                \texttt{\tiny{[GM$\rangle$P2]}}
                \texttt{ПРЕДЛОЖЕНИЕ: \{'A67', 'C70', 'C61'\}} \\
            }
        }
    }
    & \\ \\

    \theutterance \stepcounter{utterance}  
    & & & \multicolumn{4}{p{0.6\linewidth}}{
        \cellcolor[rgb]{0.9,0.9,0.9}{
            \makecell[{{p{\linewidth}}}]{
                \texttt{\tiny{[GM$\langle$P2]}}
                \texttt{АРГУМЕНТ: \{'Предложение содержит предметы с затратами: 907 + 698 + 119 = 1724. Важность: 822 + 215 + 668 = 1605. Это хороший баланс, но возможно можно улучшить набор.'\}} \\
            }
        }
    }
     \\ \\

    \theutterance \stepcounter{utterance}  
    & & \multicolumn{4}{p{0.6\linewidth}}{
        \cellcolor[rgb]{0.9,0.9,0.9}{
            \makecell[{{p{\linewidth}}}]{
                \texttt{\tiny{[P1$\langle$GM]}}
                \texttt{АРГУМЕНТ: \{'Предложение содержит предметы с затратами: 907 + 698 + 119 = 1724. Важность: 822 + 215 + 668 = 1605. Это хороший баланс, но возможно можно улучшить набор.'\}} \\
            }
        }
    }
    & \\ \\

    \theutterance \stepcounter{utterance}  
    & \multicolumn{4}{p{0.6\linewidth}}{
        \cellcolor[rgb]{0.9,0.9,0.9}{
            \makecell[{{p{\linewidth}}}]{
                \texttt{\tiny{[P1$\rangle$GM]}}
                \texttt{ПРЕДЛОЖЕНИЕ: \{'C16', 'A43', 'C73'\}} \\
            }
        }
    }
    & & \\ \\

    \theutterance \stepcounter{utterance}  
    & & \multicolumn{4}{p{0.6\linewidth}}{
        \cellcolor[rgb]{0.9,0.9,0.9}{
            \makecell[{{p{\linewidth}}}]{
                \texttt{\tiny{[GM$\rangle$P2]}}
                \texttt{ПРЕДЛОЖЕНИЕ: \{'C16', 'A43', 'C73'\}} \\
            }
        }
    }
    & \\ \\

    \theutterance \stepcounter{utterance}  
    & & & \multicolumn{4}{p{0.6\linewidth}}{
        \cellcolor[rgb]{0.9,0.9,0.9}{
            \makecell[{{p{\linewidth}}}]{
                \texttt{\tiny{[GM$\langle$P2]}}
                \texttt{АРГУМЕНТ: \{'Суммарные затраты: 866 + 868 + 805 = 2539. Важность: 856 + 868 + 508 = 2232. Этот набор может дать большую пользу, но предел ещё позволяет оптимизировать.'\}} \\
            }
        }
    }
     \\ \\

    \theutterance \stepcounter{utterance}  
    & & \multicolumn{4}{p{0.6\linewidth}}{
        \cellcolor[rgb]{0.9,0.9,0.9}{
            \makecell[{{p{\linewidth}}}]{
                \texttt{\tiny{[P1$\langle$GM]}}
                \texttt{АРГУМЕНТ: \{'Суммарные затраты: 866 + 868 + 805 = 2539. Важность: 856 + 868 + 508 = 2232. Этот набор может дать большую пользу, но предел ещё позволяет оптимизировать.'\}} \\
            }
        }
    }
    & \\ \\

    \theutterance \stepcounter{utterance}  
    & \multicolumn{4}{p{0.6\linewidth}}{
        \cellcolor[rgb]{0.9,0.9,0.9}{
            \makecell[{{p{\linewidth}}}]{
                \texttt{\tiny{[P1$\rangle$GM]}}
                \texttt{ПРЕДЛОЖЕНИЕ: \{'C16', 'A62', 'B42'\}} \\
            }
        }
    }
    & & \\ \\

    \theutterance \stepcounter{utterance}  
    & & \multicolumn{4}{p{0.6\linewidth}}{
        \cellcolor[rgb]{0.9,0.9,0.9}{
            \makecell[{{p{\linewidth}}}]{
                \texttt{\tiny{[GM$\rangle$P2]}}
                \texttt{ПРЕДЛОЖЕНИЕ: \{'C16', 'A62', 'B42'\}} \\
            }
        }
    }
    & \\ \\

    \theutterance \stepcounter{utterance}  
    & & & \multicolumn{4}{p{0.6\linewidth}}{
        \cellcolor[rgb]{0.9,0.9,0.9}{
            \makecell[{{p{\linewidth}}}]{
                \texttt{\tiny{[GM$\langle$P2]}}
                \texttt{АРГУМЕНТ: \{'Затраты: 866 + 41 + 809 = 1716. Важность: 856 + 65 + 783 = 1604. Хорошие показатели, но можно добавить предметы с высокой важностью, не превышая предел.'\}} \\
            }
        }
    }
     \\ \\

    \theutterance \stepcounter{utterance}  
    & & \multicolumn{4}{p{0.6\linewidth}}{
        \cellcolor[rgb]{0.9,0.9,0.9}{
            \makecell[{{p{\linewidth}}}]{
                \texttt{\tiny{[P1$\langle$GM]}}
                \texttt{АРГУМЕНТ: \{'Затраты: 866 + 41 + 809 = 1716. Важность: 856 + 65 + 783 = 1604. Хорошие показатели, но можно добавить предметы с высокой важностью, не превышая предел.'\}} \\
            }
        }
    }
    & \\ \\

    \theutterance \stepcounter{utterance}  
    & \multicolumn{4}{p{0.6\linewidth}}{
        \cellcolor[rgb]{0.9,0.9,0.9}{
            \makecell[{{p{\linewidth}}}]{
                \texttt{\tiny{[P1$\rangle$GM]}}
                \texttt{ПРЕДЛОЖЕНИЕ: \{'C16', 'A43', 'C70', 'C42'\}} \\
            }
        }
    }
    & & \\ \\

    \theutterance \stepcounter{utterance}  
    & & \multicolumn{4}{p{0.6\linewidth}}{
        \cellcolor[rgb]{0.9,0.9,0.9}{
            \makecell[{{p{\linewidth}}}]{
                \texttt{\tiny{[GM$\rangle$P2]}}
                \texttt{ПРЕДЛОЖЕНИЕ: \{'C16', 'A43', 'C70', 'C42'\}} \\
            }
        }
    }
    & \\ \\

    \theutterance \stepcounter{utterance}  
    & & & \multicolumn{4}{p{0.6\linewidth}}{
        \cellcolor[rgb]{0.9,0.9,0.9}{
            \makecell[{{p{\linewidth}}}]{
                \texttt{\tiny{[GM$\langle$P2]}}
                \texttt{АРГУМЕНТ: \{'Затраты: 866 + 868 + 698 + 172 = 2604. Важность: 856 + 868 + 215 + 808 = 2747. Набор даёт высокую важность и стоит в пределах.'\}} \\
            }
        }
    }
     \\ \\

    \theutterance \stepcounter{utterance}  
    & & \multicolumn{4}{p{0.6\linewidth}}{
        \cellcolor[rgb]{0.9,0.9,0.9}{
            \makecell[{{p{\linewidth}}}]{
                \texttt{\tiny{[P1$\langle$GM]}}
                \texttt{АРГУМЕНТ: \{'Затраты: 866 + 868 + 698 + 172 = 2604. Важность: 856 + 868 + 215 + 808 = 2747. Набор даёт высокую важность и стоит в пределах.'\}} \\
            }
        }
    }
    & \\ \\

    \theutterance \stepcounter{utterance}  
    & \multicolumn{4}{p{0.6\linewidth}}{
        \cellcolor[rgb]{0.9,0.9,0.9}{
            \makecell[{{p{\linewidth}}}]{
                \texttt{\tiny{[P1$\rangle$GM]}}
                \texttt{СОГЛАСЕН: \{'C16', 'A43', 'C70', 'C42'\}} \\
            }
        }
    }
    & & \\ \\

    \theutterance \stepcounter{utterance}  
    & & \multicolumn{4}{p{0.6\linewidth}}{
        \cellcolor[rgb]{0.9,0.9,0.9}{
            \makecell[{{p{\linewidth}}}]{
                \texttt{\tiny{[P1$\langle$GM]}}
                \texttt{Вы согласились с неактивным предложением. Предложения активны только в том случае, если другой игрок зарегистрировал их через ПРЕДЛОЖЕНИЕ и они не были деактивированы вами через ОТКАЗ. Попробуйте ещё раз.} \\
            }
        }
    }
    & \\ \\

    \theutterance \stepcounter{utterance}  
    & \multicolumn{4}{p{0.6\linewidth}}{
        \cellcolor[rgb]{0.9,0.9,0.9}{
            \makecell[{{p{\linewidth}}}]{
                \texttt{\tiny{[P1$\rangle$GM]}}
                \texttt{ПРЕДЛОЖЕНИЕ: \{'C10', 'B36', 'A21'\}} \\
            }
        }
    }
    & & \\ \\

    \theutterance \stepcounter{utterance}  
    & & \multicolumn{4}{p{0.6\linewidth}}{
        \cellcolor[rgb]{0.9,0.9,0.9}{
            \makecell[{{p{\linewidth}}}]{
                \texttt{\tiny{[GM$\rangle$P2]}}
                \texttt{ПРЕДЛОЖЕНИЕ: \{'C10', 'B36', 'A21'\}} \\
            }
        }
    }
    & \\ \\

    \theutterance \stepcounter{utterance}  
    & & & \multicolumn{4}{p{0.6\linewidth}}{
        \cellcolor[rgb]{0.9,0.9,0.9}{
            \makecell[{{p{\linewidth}}}]{
                \texttt{\tiny{[GM$\langle$P2]}}
                \texttt{АРГУМЕНТ: \{'Затраты: 601 + 432 + 762 = 1795. Важность: 780 + 461 + 30 = 1271. Низкая важность для моих целей, параметр затрат оптимально не использован.'\}} \\
            }
        }
    }
     \\ \\

    \theutterance \stepcounter{utterance}  
    & & \multicolumn{4}{p{0.6\linewidth}}{
        \cellcolor[rgb]{0.9,0.9,0.9}{
            \makecell[{{p{\linewidth}}}]{
                \texttt{\tiny{[P1$\langle$GM]}}
                \texttt{АРГУМЕНТ: \{'Затраты: 601 + 432 + 762 = 1795. Важность: 780 + 461 + 30 = 1271. Низкая важность для моих целей, параметр затрат оптимально не использован.'\}} \\
            }
        }
    }
    & \\ \\

    \theutterance \stepcounter{utterance}  
    & \multicolumn{4}{p{0.6\linewidth}}{
        \cellcolor[rgb]{0.9,0.9,0.9}{
            \makecell[{{p{\linewidth}}}]{
                \texttt{\tiny{[P1$\rangle$GM]}}
                \texttt{ПРЕДЛОЖЕНИЕ: \{'C67', 'B40', 'C39'\}} \\
            }
        }
    }
    & & \\ \\

    \theutterance \stepcounter{utterance}  
    & & \multicolumn{4}{p{0.6\linewidth}}{
        \cellcolor[rgb]{0.9,0.9,0.9}{
            \makecell[{{p{\linewidth}}}]{
                \texttt{\tiny{[GM$\rangle$P2]}}
                \texttt{ПРЕДЛОЖЕНИЕ: \{'C67', 'B40', 'C39'\}} \\
            }
        }
    }
    & \\ \\

    \theutterance \stepcounter{utterance}  
    & & & \multicolumn{4}{p{0.6\linewidth}}{
        \cellcolor[rgb]{0.9,0.9,0.9}{
            \makecell[{{p{\linewidth}}}]{
                \texttt{\tiny{[GM$\langle$P2]}}
                \texttt{АРГУМЕНТ: \{'Затраты: 462 + 248 + 163 = 873. Важность: 389 + 97 + 500 = 986. Набор мал по сумме важности, но безопасен в пределах.'\}} \\
            }
        }
    }
     \\ \\

    \theutterance \stepcounter{utterance}  
    & & \multicolumn{4}{p{0.6\linewidth}}{
        \cellcolor[rgb]{0.9,0.9,0.9}{
            \makecell[{{p{\linewidth}}}]{
                \texttt{\tiny{[P1$\langle$GM]}}
                \texttt{АРГУМЕНТ: \{'Затраты: 462 + 248 + 163 = 873. Важность: 389 + 97 + 500 = 986. Набор мал по сумме важности, но безопасен в пределах.'\}} \\
            }
        }
    }
    & \\ \\

    \theutterance \stepcounter{utterance}  
    & \multicolumn{4}{p{0.6\linewidth}}{
        \cellcolor[rgb]{0.9,0.9,0.9}{
            \makecell[{{p{\linewidth}}}]{
                \texttt{\tiny{[P1$\rangle$GM]}}
                \texttt{ПРЕДЛОЖЕНИЕ: \{'C55', 'A18', 'B19'\}} \\
            }
        }
    }
    & & \\ \\

    \theutterance \stepcounter{utterance}  
    & & \multicolumn{4}{p{0.6\linewidth}}{
        \cellcolor[rgb]{0.9,0.9,0.9}{
            \makecell[{{p{\linewidth}}}]{
                \texttt{\tiny{[GM$\rangle$P2]}}
                \texttt{ПРЕДЛОЖЕНИЕ: \{'C55', 'A18', 'B19'\}} \\
            }
        }
    }
    & \\ \\

    \theutterance \stepcounter{utterance}  
    & & & \multicolumn{4}{p{0.6\linewidth}}{
        \cellcolor[rgb]{0.9,0.9,0.9}{
            \makecell[{{p{\linewidth}}}]{
                \texttt{\tiny{[GM$\langle$P2]}}
                \texttt{АРГУМЕНТ: \{'Затраты: 118 + 166 + 193 = 477. Важность: 583 + 484 + 262 = 1329. Высокая ценность при минимальных затратах — перспективное решение.'\}} \\
            }
        }
    }
     \\ \\

    \theutterance \stepcounter{utterance}  
    & & \multicolumn{4}{p{0.6\linewidth}}{
        \cellcolor[rgb]{0.9,0.9,0.9}{
            \makecell[{{p{\linewidth}}}]{
                \texttt{\tiny{[P1$\langle$GM]}}
                \texttt{АРГУМЕНТ: \{'Затраты: 118 + 166 + 193 = 477. Важность: 583 + 484 + 262 = 1329. Высокая ценность при минимальных затратах — перспективное решение.'\}} \\
            }
        }
    }
    & \\ \\

    \theutterance \stepcounter{utterance}  
    & \multicolumn{4}{p{0.6\linewidth}}{
        \cellcolor[rgb]{0.9,0.9,0.9}{
            \makecell[{{p{\linewidth}}}]{
                \texttt{\tiny{[P1$\rangle$GM]}}
                \texttt{СОГЛАСЕН: \{'C55', 'A18', 'B19'\}} \\
            }
        }
    }
    & & \\ \\

    \theutterance \stepcounter{utterance}  
    & & & \multicolumn{2}{p{0.3\linewidth}}{
        \cellcolor[rgb]{0.9,0.9,0.9}{
            \makecell[{{p{\linewidth}}}]{
                \texttt{\tiny{[GM$|$GM]}}
                \texttt{Maximum rounds exceeded: 10 >= 10} \\
            }
        }
    }
    & & \\ \\

    \theutterance \stepcounter{utterance}  
    & & & \multicolumn{2}{p{0.3\linewidth}}{
        \cellcolor[rgb]{0.9,0.9,0.9}{
            \makecell[{{p{\linewidth}}}]{
                \texttt{\tiny{[GM$|$GM]}}
                \texttt{end game} \\
            }
        }
    }
    & & \\ \\

\end{supertabular}
}

\end{document}
